% Extensión aditiva. Canon leído: CURRENT.json rev. 2026-07-22.2.
% Autores del proyecto: Oumar Haidara Fall y Rubén Ramos Balsa.

\chapter{Lines of magnitude and dimensional realization of the HMT kernel}
\label{chap:funtor-dimensional-hmt}

The transition from HMT to Dimensional Mechanics does not consist in adjoining a unit symbol to a real number. The arithmetic core produces characters, cocycles, ratios, and orders of a decimal chart. By contrast, a physical magnitude is an element of a line endowed with dimensional type. The distinction is structural: a change of units modifies the coordinates of a magnitude, but not the magnitude itself. On the one hand, this section constructs the ambient lattice required to realize mass, length, duration, and charge independently, and determines that its minimum rank is four. On the other hand, it constructs a monoidal functor of trajectories with signature that takes values in a portion of that lattice. The minimality of the lattice is not attributed to the functor.

The construction has two consequences. First, it makes it possible to realize action, duration, length, charge, and mass in a common lattice without introducing the International System into the generative formulas; assigning charge to states additionally requires a specific integral character. Second, it proves exactly which part of an absolute scale cannot follow from dimensionless scalars. This impossibility does not affect ratios, mass characters, or constitutive identities; it fixes the correct mathematical type of their outputs. The distinction among dimension, numerical coordinate, and choice of unit is that of the classical theory of dimensional analysis \cite{Buckingham1914,Bridgman1931}; the structure of lines that follows makes this distinction explicit in the HMT domain\@.

\section{Lattice of dimensions and oriented lines}

Let
\[
 \Lambda_{\rm D}
 =\mathbb Z\langle \mathbf M,\mathbf L,\mathbf T,\mathbf Q\rangle
 \cong\mathbb Z^4
\]
the lattice of dimensions, written in terms of mass, length, duration, and
charge. For \(d=(d_M,d_L,d_T,d_Q)\in\Lambda_{\rm D}\), we denote by \(\mathcal L_d\) an oriented real line
and by \(\mathcal L_d^+\) its positive ray. The tensor product realizes addition
of dimensions,
\[
 \mathcal L_d\otimes\mathcal L_e\simeq\mathcal L_{d+e},
 \qquad
 \mathcal L_d^\vee\simeq\mathcal L_{-d}.
\]
The collection of these lines, with positive isomorphisms and tensor
product, forms a symmetric Picard groupoid. A dimensionless scalar is
an element of \(\mathcal L_0\simeq\mathbb R\); a quantity of dimension
\(d\) is an element of \(\mathcal L_d\).

We introduce four dimensions adapted to HMT:
\[
 \mathbf S=\mathbf M+2\mathbf L-\mathbf T,
 \qquad
 \mathbf C=\mathbf L-\mathbf T,
 \qquad
 \mathbf T=\mathbf T,
 \qquad
 \mathbf Q=\mathbf Q.
\]
They represent, respectively, action, velocity, duration, and charge. In the
basis \((\mathbf M,\mathbf L,\mathbf T,\mathbf Q)\), their columns form the
matrix
\[
 B_{\rm D}=
 \begin{pmatrix}
 1&0&0&0\\
 2&1&0&0\\
 -1&-1&1&0\\
 0&0&0&1
 \end{pmatrix}.
\]

\begin{theorem}[Adapted base dimension for HMT]
\label{thm:base-dimensional-hmt}
The dimensions
\[
 \mathbf S,\qquad\mathbf C,\qquad\mathbf T,\qquad\mathbf Q
\]
constitute a unimodular basis of
\(\Lambda_{\rm D}\). In particular, every mechanical,
electromagnetic, or gravitational dimension generated by
\((\mathbf M,\mathbf L,\mathbf T,\mathbf Q)\) has a unique expression with
integer exponents in the four lines
\(\mathcal L_S,\mathcal L_C,\mathcal L_T,\mathcal L_Q\).
\end{theorem}

\begin{proof}
The matrix \(B_{\rm D}\) is block triangular and
\(\det B_{\rm D}=1\). It therefore belongs to
\(\operatorname{GL}_4(\mathbb Z)\). Its integral inverse provides
\[
 \mathbf L=\mathbf C+\mathbf T,
 \qquad
 \mathbf M=\mathbf S-2\mathbf C-\mathbf T.
\]
The four original dimensions are thereby uniquely recovered.
\end{proof}

Among the identities that we will use are
\begin{align}
 \mathbf E&=\mathbf S-\mathbf T,
 &\mathbf Z&=\mathbf S-2\mathbf Q,\label{eq:dim-EZ}\\
 \boldsymbol\mu&=\mathbf S-\mathbf C-2\mathbf Q,
 &\boldsymbol\varepsilon&=-\mathbf S-\mathbf C+2\mathbf Q,
 \label{eq:dim-mueps}\\
 \mathbf G&=-\mathbf S+5\mathbf C+2\mathbf T,
 &\boldsymbol\kappa&=-\mathbf S+\mathbf C+2\mathbf T,
 \label{eq:dim-Gkappa}\\
 \boldsymbol\Lambda&=-2\mathbf C-2\mathbf T.
 \label{eq:dim-Lambda}
\end{align}
Here \(\mathbf E\) is energy, \(\mathbf Z\) impedance,
\(\boldsymbol\mu\) permeability, \(\boldsymbol\varepsilon\)
permittivity, \(\mathbf G\) the dimension of the gravitational constant,
\(\boldsymbol\kappa=\mathbf G-4\mathbf C\), and
\(\boldsymbol\Lambda=-2\mathbf L\).

\section{Rescaling group and impossibility theorem}

The group of unit base change transformations is
\[
 \mathcal G_{\rm D}
 :=\operatorname{Hom}(\Lambda_{\rm D},\mathbb R_{>0})
 \cong(\mathbb R_{>0})^4.
\]
For \(g\in\mathcal G_{\rm D}\) and \(d\in\Lambda_{\rm D}\), write
\(g^d\) for the corresponding character. If \(u_d\) is a positive basis
of \(\mathcal L_d\), the chart change
\(u_d\mapsto g^d u_d\) transforms the coordinate of a fixed quantity according to
\[
 x\,u_d=(g^{-d}x)(g^du_d).
\]
The magnitude does not change; its numerical representation changes.

\begin{theorem}[Impossibility of an Absolute Scale from Dimensionless Data]
\label{thm:no-go-dimensional-rango-cuatro}
Let \(X\) be a space on which \(\mathcal G_{\rm D}\) acts
trivially, and let \(d\ne0\). Every map
\(F:X\to\mathcal L_d\) that is natural with respect to all changes of
units and that receives no additional dimensional basis is identically
zero. By contrast, every quotient of two nonzero sections of the same line,
and every product whose total dimension is zero, is invariant under
\(\mathcal G_{\rm D}\).
\end{theorem}

\begin{proof}
Naturality and the trivial action on \(X\) imply
\[
 F(x)=F(gx)=g^dF(x)
 \qquad(g\in\mathcal G_{\rm D}).
\]
Since \(d\ne0\), there exists \(g\) with \(g^d\ne1\); hence
\((g^d-1)F(x)=0\), and then \(F(x)=0\). For two sections of dimension
\(d\), the factors \(g^d\) cancel in the quotient. The same cancellation
occurs in every monomial of total dimension zero.
\end{proof}

The theorem does not establish that a discrete theory is incapable of constructing an
internal scale. It establishes that this scale must appear as a section of
a line or, equivalently, that the theory must extend its domain by the
torsor of dimensional frames. It is precisely this extension that defines
Dimensional Mechanics.

\begin{corollary}[Minimum rank]
Every system of lines that permits mass, length, duration, and charge to be realized
independently has scale rank at least four. The basis
\((\mathbf S,\mathbf C,\mathbf T,\mathbf Q)\) attains the bound; it is therefore
minimal.
\end{corollary}

\begin{proof}
The four original dimensions are linearly independent in
\(\Lambda_{\rm D}\). Three generators cannot span a rank-four
lattice. The \cref{thm:base-dimensional-hmt} provides four generators that
do span it.
\end{proof}

\section{Internal dimensional space: degrees, charts, and HMT–MD transduction operators}

An internal framework HMT--MD is a quadruple of positive bases
\[
 \mathfrak u=(u_S,u_C,u_T,u_Q)
 \in\operatorname{Fr}_{\rm D},
\]
where \(\operatorname{Fr}_{\rm D}\) is a principal torsor under
\(\mathcal G_{\rm D}\). The derived bases are defined tensorially:
\begin{align*}
 u_L&=u_Cu_T,&
 u_M&=u_Su_C^{-2}u_T^{-1},&
 u_E&=u_Su_T^{-1},\\
 u_Z&=u_Su_Q^{-2},&
 u_\mu&=u_Su_C^{-1}u_Q^{-2},&
 u_\varepsilon&=u_S^{-1}u_C^{-1}u_Q^2,\\
 u_G&=u_S^{-1}u_C^5u_T^2,&
 u_\kappa&=u_S^{-1}u_Cu_T^2,&
 u_\Lambda&=u_C^{-2}u_T^{-2}.
\end{align*}
Juxtaposition here denotes the tensor product of bases. An external chart
is another frame \(\mathfrak u_{\rm ext}\) together with a positive isomorphism
between the two systems of lines. Nothing requires that chart to be SI; the
choice of SI belongs to the metrological publication of a magnitude.

\subsection{Action line}

The determinantal reader local begins from
\[
 G_H=\operatorname{coker}(H_4),
 \qquad |G_H|=16,
 \qquad \Sigma_H=\varphi\alpha_{\rm HMT}^{16},
\]
and proves
\[
 10^{-34}<\Sigma_H<10^{-33},
 \qquad \operatorname{ord}_{10}(\Sigma_H)=-34.
\]
The post-return regional realization provides the dimensionless angular coordinate.
\[
 \etaRet
 =\frac{\Cdeg}{2}-\alpha_{\angle,{\rm deg}}
 =1.0545718169150390611336905847242997\ldots.
\]
We define the internal action reduced section by
\begin{equation}
 \hbar_{\rm int}
 :=\etaRet\,10^{-34}u_S,
 \qquad
 h_{\rm int}:=2\pi\hbar_{\rm int}.
 \label{eq:accion-interna}
\end{equation}
The exponent is an output of the reader; \(u_S\) is the basis of a line.
The composition is covariant under \(\mathcal G_{\rm D}\) and contains no
identification with \(\mathrm{J\,s}\).

\subsection{Periodicity of route edge and resulting constitutive velocity}

Let \(\mathsf P_{\rm TPK}\) be the free category of trajectories of the graph of
enriched states. The grading by length
\[
 \tau(\gamma)=|\gamma|,
 \qquad
 \tau(\gamma\eta)=\tau(\gamma)+\tau(\eta)
\]
define the combinatorial clock of the protocol. Its implementation is
\[
 T_{\rm int}(\gamma)=\tau(\gamma)u_T.
\]
This duration counts transitions of the fixed graph. A subdivision of the graph
defines another clock and must be accompanied by its refinement morphism.

The channels \(q_\pm\) generate the symmetric and antisymmetric functions
\[
 \mathcal E=
 \log\!\bigl[(1-q_+^{90})(1-q_-^{90})\bigr]
 -\log\!\bigl[(1-q_+^{120})(1-q_-^{120})\bigr],
\]
\[
 \mathcal B=
 \log\!\frac{1-q_-^{90}}{1-q_+^{90}}
 -\log\!\frac{1-q_-^{120}}{1-q_+^{120}}.
\]
These determine
the positive coefficients
\[
 \widehat c=e^{-\mathcal E},
 \qquad \widehat Z=e^{\mathcal B}.
\]
Its internal realizations are
\begin{equation}
 c_{\rm int}=\widehat c\,u_C,
 \qquad Z_{\rm int}=\widehat Z\,u_Z.
 \label{eq:cZ-internos}
\end{equation}
For a constant channel along a path one obtains the length
\begin{equation}
 L_{\rm int}(\gamma)
 =\widehat c\,|\gamma|\,u_L.
 \label{eq:longitud-ruta}
\end{equation}
If the channel depends on the edge, the covariant formula is the sum
\(\sum_{e\in\gamma}\widehat c(e)u_L\); its additivity under concatenation is
preserved.

\subsection{Elementary charge scale and conditional character}

Since \(u_Z=u_Su_Q^{-2}\), the quotient
\(h_{\rm int}/Z_{\rm int}\) belongs to \(\mathcal L_Q^{\otimes2}\).
The positive orientation of the lines selects a unique positive root.
We define
\begin{equation}
 e_{\rm int}
 :=\left(\frac{2\alpha_{\rm HMT}h_{\rm int}}
 {Z_{\rm int}}\right)^{1/2}\in\mathcal L_Q^+.
 \label{eq:carga-interna}
\end{equation}

\begin{proposition}[Internal Closure of the Fine Structure Relation]
The section \eqref{eq:carga-interna} is the unique positive section that
satisfies
\[
 \alpha_{\rm HMT}
 =\frac{Z_{\rm int}e_{\rm int}^{\,2}}{2h_{\rm int}}.
\]
The equality is invariant under the full scaling group.
\end{proposition}

\begin{proof}
Existence follows from the positivity of
\(\alpha_{\rm HMT},h_{\rm int},Z_{\rm int}\). An oriented real line
has a unique positive root of a positive section of its square. The
dimension of the right-hand side is
\(\mathbf Z+2\mathbf Q-\mathbf S=0\), so all rescaling factors
cancel.
\end{proof}

This proposition does not count as a second independent determination of
\(\alpha\): within the system of lines, it realizes the constitutive
compatibility among action, impedance, and an elementary charge scale. It does not
construct a charge on the states. This additionally requires an
integer character
\[
 n_Q:\operatorname{Ob}(\mathsf P_{\rm TPK}^{\rm sig})\longrightarrow
 \mathbb Z;
 \qquad
 Q(x)=n_Q(x)e_{\rm int}.
\]
The existence and selection of \(n_Q\) are additional hypotheses. In
particular, \(n_Q\) is not a component of the trajectory functor
constructed below. The elementary scale and the classification of states
are distinct maps.

\section{Monoidal path functor with signature}

Let \(\mathsf P_{\rm TPK}^{\rm sig}\) be the category of enriched
paths decorated with a signature coordinate
\(v_x\in\mathbb Z^5\) at each object. For
\(\gamma:x\to y\), we set
\[
 \Delta v(\gamma)=v_y-v_x.
\]
Composition satisfies
\(\Delta v(\gamma\eta)=\Delta v(\gamma)+\Delta v(\eta)\).
Let
\[
 \chi_{\rm mass}(v)=\exp\langle\lambda_{\rm mass},v\rangle
\]
the positive mass character. On the route quotient, we fix the
action law determined by five data: the normalized trace of the distinct
emissions, a positive additive cocycle, the change defect
\(\Delta_4P_3\), the residual capacity \(\lambda=1-\delta\), and the local sum
of edge costs. In this law,
\[
 \mathcal A_*(\gamma)
 =N_{\rm giro}(\gamma)+\lambda_*N_{\rm hoja}(\gamma),
 \qquad
 \lambda_*=1-\frac5{468}-\frac3{11}\Delta_4.
\]
The following argument requires only additivity; it also applies to the
functional \(I_\lambda\) with any \(\lambda>0\).

Consider the monoid
\[
 \mathfrak M_{\rm D}
 =\mathbb N\times\mathbb R_{\ge0}\times
  \mathbb R_{\ge0}\times\mathbb R_{>0}
\]
with product
\[
 (t,\ell,s,r)\odot(t',\ell',s',r')
 =(t+t',\ell+\ell',s+s',rr').
\]

\begin{theorem}[Monoidal path functor with signature]
\label{thm:funtor-dimensional-minimo}
The map
\[
 \mathfrak D_{\rm HMT}(\gamma)
 =\bigl(
 |\gamma|,
 \widehat c|\gamma|,
 \mathcal A_*(\gamma),
 \chi_{\rm mass}(\Delta v(\gamma))
 \bigr)
\]
define a monoidal functor
\[
 \mathfrak D_{\rm HMT}:
 \mathsf P_{\rm TPK}^{\rm sig}\longrightarrow\mathfrak M_{\rm D}.
\]
After choosing an internal frame, its three additive components are realized as
\[
 T_{\rm int}(\gamma)=|\gamma|u_T,
 \quad
 L_{\rm int}(\gamma)=\widehat c|\gamma|u_L,
 \quad
 S_{\rm int}(\gamma)=\mathcal A_*(\gamma)\hbar_{\rm int},
\]
and the fourth transports dimensionless mass ratios. The functor does not contain
a charge coordinate and does not generate the ambient dimensional lattice.
\end{theorem}

\begin{proof}
The length of a trajectory, the number of turns, and the number of changes of
sheet are additive under compatible concatenation. By hypothesis,
\(\mathcal A_*\) is additive. The difference of signatures is additive, and the mass
character transforms sums into products. Each component therefore respects
the law \(\odot\), and the empty route is sent to
\((0,0,0,1)\). The linear realization preserves sums because it multiplies
by fixed sections of the corresponding lines.
\end{proof}

\begin{remark}[Ambiental lattice and scope of the functor]
The minimum-rank corollary concerns a system capable of varying
mass, length, duration, and charge independently. The preceding
functor realizes duration, length, and action, and transports a mass ratio; it does not use
\(u_Q\) or prove a rank bound. The subsequent assignment of charges
remains conditional upon the existence of \(n_Q\). In the dimensional lattice,
the types of its components are
\[
 \mathbf T,\qquad \mathbf L=\mathbf C+\mathbf T,\qquad
 \mathbf S,\qquad \mathbf0,
\]
so that its dimensional image is contained in
\(\langle\mathbf S,\mathbf C,\mathbf T\rangle_{\mathbb Z}\), of rank three.
The rescaling \(u_Q\mapsto\lambda u_Q\), with the other bases fixed, leaves
all components of the functor invariant; the latter does not faithfully realize the
fourth ambient direction.
\end{remark}

If the bidirectional electronic section is assigned to each object as
the origin of the mass chart,
\[
 m(x)=m_e^{\beta}\chi_{\rm mass}(v_x-v_e),
\]
then every morphism \(\gamma:x\to y\) satisfies exactly
\[
 \frac{m(y)}{m(x)}
 =\chi_{\rm mass}(\Delta v(\gamma)).
\]
The theorem closes dimensional propagation once the signature is given. It
does not change the status of the preceding map associating a species with
a trajectory and its signature.

\section{Basal stage and bidirectional closure of the electronic scale}

The prevailing electronic invariant is the positive scalar.
\[
 \mathfrak e_0=
 \frac{\sqrt3}{4}
 \left(e^{\varphi/\pi^2}-22\alpha_{\rm HMT}^3\right)
 (1+15\Delta_4).
\]
Its numerical value is
\(0.510998922488066\ldots\). This is the basal stage of the genealogy, not its
governing output. In the internal dimensional framework, it defines an energy section
and a mass section related by the velocity:
\begin{equation}
 E_e^{\rm int}=\mathfrak e_0u_E,
 \qquad
 m_e^{\rm int}=\frac{E_e^{\rm int}}{c_{\rm int}^{2}}
 =\mathfrak e_0\widehat c^{-2}u_M.
 \label{eq:electron-line-valued}
\end{equation}
The central route and its two bidirectional readers determine, without consulting
the external mass,
\[
 \kappa_e=80-54\alpha_{\rm HMT}+6\Delta_4,
 \qquad R_{\rm act}=\frac{H_5}{\etaRet}.
\]
The current electronic section is, therefore,
\begin{equation}
 \begin{aligned}
 E_e^{\beta}&=E_e^{(0)}R_{\rm act}^{\kappa_e},
 &m_e^{\beta}&=\frac{E_e^{\beta}}{c_{\rm int}^{2}},\\
 E_e^\beta
 &=0.5109989506900008086082571648\ldots\ {\rm MeV},
 &u_E&=1\,{\rm MeV}.
 \end{aligned}
 \label{eq:electron-line-valued-beta}
\end{equation}
Thus, the basal formula retains its derivative role, and bidirectional
closure governs all subsequent propagation through characters and towers.
In a natural-units chart in which the coordinate of \(c_{\rm int}\) is fixed
at one, both numerical coordinates coincide; their lines are not silently
identified.

\begin{proposition}[Typed separation of the energy base and the clock quantum]
\label{prop:no-identificacion-base-reloj-rev8}
Escribamos \(t_0=\tau_0u_T\), with \(\tau_0>0\), and definamos
\[
 \varepsilon_{\rm clk}
 :=\hbar_{\rm int}\frac{2\pi}{108t_0},
 \qquad
 \kappa_{\rm clk}(\tau_0)
 :=\frac{2\pi\etaRet\,10^{-34}}{108\tau_0}.
\]
Then
\[
 \varepsilon_{\rm clk}=\kappa_{\rm clk}(\tau_0)u_E,
 \qquad
 m_0:=\frac{\varepsilon_{\rm clk}}{c_{\rm int}^2}
 =\kappa_{\rm clk}(\tau_0)\widehat c^{-2}u_M,
\]
and the current electronic section satisfies
\[
 E_e^{\rm int}
 =\frac{\mathfrak e_0}{\kappa_{\rm clk}(\tau_0)}
  \varepsilon_{\rm clk}.
\]
Therefore, \(E_e^{\rm int}=\mathfrak e_0\varepsilon_{\rm clk}\) if and only if
\(\kappa_{\rm clk}(\tau_0)=1\). Adopting \(\varepsilon_{\rm clk}\) as the basis
while retaining the same coordinate \(\mathfrak e_0\) is a new choice of
frame; it is not a correction of the electronic invariant, nor does it follow
from order-108 periodicity. In particular, the normalization used
later, \(t_0=u_T\), fixes \(\tau_0=1\), not
\(\kappa_{\rm clk}=1\).
\end{proposition}

\begin{proof}
One substitutes \(\hbar_{\rm int}=\etaRet\,10^{-34}u_S\), \(t_0=\tau_0u_T\), \(u_E=u_Su_T^{-1}\), and
\(c_{\rm int}=\widehat c,u_C\). The final identity follows by comparison with
\cref{eq:electron-line-valued}. All equalities occur between sections of the same line; none
consults an external unit or mass.
\end{proof}

\begin{proposition}[Mass scale propagation]
For each declared signature \(v\in\mathbb Z^5\),
\[
 m_v^{\rm int}:=m_e^{\rm int}\chi_{\rm mass}(v)
 \]
is a positive section of \(\mathcal L_M\), and
\[
 \frac{m_{v+w}^{\rm int}}{m_v^{\rm int}}
 =\chi_{\rm mass}(w).
\]
The construction is covariant under rescaling and does not require an external unit of energy or mass.
\end{proposition}

\begin{proof}
The character is dimensionless and positive, so it does not modify the line of
\(m_e^{\rm int}\). The quotient identity expresses the multiplicativity of
\(\chi_{\rm mass}\). Under a change of frame, the numerator and denominator
acquire the same weight \(\mathbf M\), which cancels.
\end{proof}

A subsequent metrological chart may send \(u_E\) to an energy basis
and \(u_M\) to the coherent basis obtained by dividing by the square
of the standard velocity. Only that chart authorizes the section to be
expressed by an external unit name.

\section{Constitutive as isomorphisms between lines}

The internal functions satisfy
\[
 \widehat\mu=e^{\mathcal B+\mathcal E},
 \qquad
 \widehat\varepsilon=e^{-\mathcal B+\mathcal E},
 \qquad
 \widehat Z=e^{\mathcal B},
 \qquad
 \widehat c=e^{-\mathcal E}.
\]
We define the sections
\[
 \mu_{\rm int}=\widehat\mu u_\mu,
 \qquad
 \varepsilon_{\rm int}=\widehat\varepsilon u_\varepsilon.
\]

\begin{theorem}[Covariant Linear Constitutive Relations]
In the groupoid of oriented lines, the exact identities hold.
\[
 \sqrt{\frac{\mu_{\rm int}}{\varepsilon_{\rm int}}}
 =Z_{\rm int},
 \qquad
 \frac1{\sqrt{\mu_{\rm int}\varepsilon_{\rm int}}}
 =c_{\rm int}.
\]
The four sections depend on two dimensionless characters and on the two
lines \(\mathcal L_Z,\mathcal L_C\); they do not constitute four independent
scales.
\end{theorem}

\begin{proof}
In coefficients,
\(\widehat\mu/\widehat\varepsilon=\widehat Z^2\) and
\(\widehat\mu\widehat\varepsilon=\widehat c^{-2}\). In dimensions,
\[
 \boldsymbol\mu-\boldsymbol\varepsilon=2\mathbf Z,
 \qquad
 \boldsymbol\mu+\boldsymbol\varepsilon=-2\mathbf C.
\]
The positive roots are therefore sections of the correct lines, and the
identities are covariant.
\end{proof}

\section{Cartan Scale and Homogeneity of the Gravitational Action}

The elementary duration \(t_0=u_T\) and the constitutive velocity determine the
length of a transition,
\[
 \ell_0=c_{\rm int}t_0=\widehat c\,u_L.
\]
This section allows for the homogenization of the Cartan affine connection:
\[
 \mathcal A_{\ell_0}
 =\begin{pmatrix}
 \omega^I{}_J&\ell_0^{-1}e^I\\
 0&0
 \end{pmatrix},
 \qquad
 \mathcal F_{\ell_0}
 =\begin{pmatrix}
 F^I{}_J&\ell_0^{-1}T_{\rm tor}^I\\
 0&0
 \end{pmatrix}.
\]
The quotient \(e/\ell_0\) is dimensionless; curvature and torsion remain in
distinct blocks of the same transport.

The dimension of \(G\) does not add a fifth generator, since
\(\mathbf G=-\mathbf S+5\mathbf C+2\mathbf T\). For every positive
dimensionless parameter \(\theta_G\), we define the covariant family
\begin{equation}
 G_{\rm int}(\theta_G)
 :=\theta_G\frac{c_{\rm int}^{5}t_0^2}{\hbar_{\rm int}}
 \in\mathcal L_G^+.
 \label{eq:G-interno}
\end{equation}
The choice \(\theta_G=1\) is the internal Planck normalization associated with
an elementary transition. By construction, it satisfies
\[
 \ell_0^2=\frac{\hbar_{\rm int}G_{\rm int}(1)}
 {c_{\rm int}^{3}}.
\]
This reference normalization is neither the circular selector nor the
holonomic selector of the physical section of \(G\).
The parameter \(\theta_G\) makes visible the freedom not fixed by the
dimensional identities. A TPK dynamics that produces a gravitational
character must determine it before any external chart.

Let
\[
 \kappa_{\rm int}=\frac{8\pi G_{\rm int}}{c_{\rm int}^{4}}.
\]
In the following functional, \(M\) is a smooth oriented manifold of
dimension four, and \(\gamma\in\mathbb R\setminus\{0\}\). If \(M\) has a
boundary, variations of compact support in the interior are imposed, or a
boundary term with conditions that cancel the boundary variation. A
spinorial matter action additionally presupposes a compatible spin
structure.
If the co-tetrad \(e^I\) takes values in \(\mathcal L_L\) and the curvature
\(F^{IJ}\) is a dimensionless two-form after the connection has been integrated,
then
\(\int\Sigma_{IJ}\wedge F^{IJ}\) has dimension \(2\mathbf L\).

\begin{theorem}[Homogeneity of the Palatini--Cartan--Holst Action]
\label{thm:homogeneidad-pch}
With explicit units, the gravitational functional of action dimension is
\[
 \boxed{
 S_{\rm PCH}
 =\frac1{2\kappa_{\rm int}c_{\rm int}}
 \int\Sigma_{IJ}\wedge
 \left(\star F^{IJ}+\gamma^{-1}F^{IJ}\right)
 -\frac1{\kappa_{\rm int}c_{\rm int}}
 \int\Lambda_{\rm int}\operatorname{vol}_e
 +S_{\rm m}.}
 \]
Both geometric terms belong to \(\mathcal L_S\). Writing them with
\(1/(2\kappa)\) without the factor \(c^{-1}\) belongs to a convention
\(c=1\) or has dimension \(\mathbf S+\mathbf C\), not action dimension.
\end{theorem}

\begin{proof}
De \eqref{eq:dim-Gkappa},
\(\boldsymbol\kappa=-\mathbf S+\mathbf C+2\mathbf T\). Moreover,
\(2\mathbf L=2\mathbf C+2\mathbf T\). therefore,
\[
 -\boldsymbol\kappa-\mathbf C+2\mathbf L
 =\mathbf S.
\]
The cosmological term has the same dimension because
\(\boldsymbol\Lambda+4\mathbf L=2\mathbf L\). If
\(-\mathbf C\) is omitted, the result is \(\mathbf S+\mathbf C\).
\end{proof}

The Cartan--Holst equation then retains its variational content
provided that the spin current is defined with the normalization consistent
with this functional. The dynamical realization must still provide a
morphism from the TPK holonomies to \((e,\omega,\Psi)\); the present
section constructs the exact dimensional codomain and removes the ambiguity of
the units in that morphism.

\section{Universality of Dimensional Realization}

We collect the results. Let \(\mathcal X_{\rm HMT}\) be the domain of enriched
states, and let the following be dimensionless data on it:
\[
 \alpha_{\rm HMT},\quad \Adeg,\quad\Cdeg,\quad
 \widehat c,\widehat Z,\quad
 \mathcal A_*,\quad\chi_{\rm mass},\quad\mathfrak e_0,
\]
built in the previous chapters over their declared domains.
Consider the extended domain
\[
 \mathcal X_{\rm MD}
 :=\mathcal X_{\rm HMT}\times\operatorname{Fr}_{\rm D}.
\]

\begin{theorem}[Universal dimensional extension]
\label{thm:extension-dimensional-universal}
There exists a unique tensorial, positive, and
\(\mathcal G_{\rm D}\)-equivariant extension that:
\begin{enumerate}
 \item realizes the edge clock in \(\mathcal L_T\);
 \item realizes \(\widehat c\) and \(\widehat Z\) in
       \(\mathcal L_C\) and \(\mathcal L_Z\);
 \item realizes the determinantal index and the post-return component
       in \(\mathcal L_S\);
 \item satisfies the exact relations among impedance, propagation
       speed, permeability, and permittivity;
 \item satisfies the internal fine-structure relation with an elementary scale
       of positive charge and, if a character \(n_Q\) is supplied, conditionally realizes
       \(Q(x)=n_Q(x)e_{\rm int}\);
 \item transports the mass character into \(\mathcal L_M\).
\end{enumerate}
All other lines of the lattice are obtained by tensor products and
duals. Two realizations differ by exactly one element of
\(\mathcal G_{\rm D}\).
\end{theorem}

\begin{proof}
The \cref{thm:base-dimensional-hmt} identifies \((\mathbf S,\mathbf C,\mathbf T,\mathbf Q)\) with a free basis of the lattice. A monoidal map on a free category is determined by its value on the generators. Positivity fixes the roots that appear in charge and the constitutive relations. The identities already proved ensure that the definitions do not depend on an alternative presentation of a dimension. The simply transitive action of \(\mathcal G_{\rm D}\) on the frames proves the final assertion.
\end{proof}

The mathematical conclusion is precise. HMT fixes the genealogy of the
internal coefficients, and Dimensional Mechanics realizes them in a minimal
ambient lattice of rank four when mass, length, duration, and charge must
vary independently. The monoidal functor of trajectories with signature
occupies part of that lattice and therefore does not inherit a claim of
minimality. Ratios and homogeneous equations are independent of every chart.
A publication in seconds, joules, coulombs, or
electronvolts is the subsequent composition with an external frame; it does not alter
any of the internal theorems and cannot replace them.
